\subsection{Mechanical subsystem}

Let $x(t)$ be the ball position measured from the top of the pedestal, 
$r$ the ball radius and $l$ the distance between the pedestal reference 
and the electromagnet. The distance between the ball and the electromagnet is
\begin{equation}
    d(t) = l - x(t) - 2r .
    \label{eq:air_gap}
\end{equation}

The magnetic force generated by the electromagnet is modelled as

\begin{equation}
    F_m(d,i) = K_m \frac{i^2(t)}{d^2(t)} ,
    \label{eq:magnetic_force}
\end{equation}

where $K_m$ is the electromagnet force constant and $i(t)$ is the coil current.
Unfortunately, $K_m$ parameter introduces a significant uncertainty in the model, since it influenes the force generated by the electomagnet,
but it is not possible to measure it without a dynamometer, so it is used the datasheet value.

Newton's law gives
\begin{equation}
    m\ddot{x}(t) = F_m(d,i) - mg .
    \label{eq:newton}
\end{equation}

Substituting \eqnref{eq:magnetic_force} and \eqnref{eq:air_gap} into 
\eqnref{eq:newton}, the nonlinear mechanical equation becomes

\begin{equation}
    \ddot{x}(t) =
    \frac{K_m}{m}
    \frac{i^2(t)}{\left(l-x(t)-2r\right)^2}
    - g .
    \label{eq:mechanical_nonlinear}
\end{equation}